% !TEX root = master.tex
\chapter{\"Uberblick und Umfang}
\label{ch:b-ueberblick}

\emph{SunShine~Tours} ist eine als Demonstrationsprojekt umgesetzte Web-Seite f\"ur einen
fiktiven Anbieter von St\"adtereisen. Die inhaltliche Ausgestaltung ist bewusst satirisch
gehalten; technisch handelt es sich um eine serverseitig mit \ac{php} erzeugte Web-Anwendung,
die ohne Datenbank auskommt und ihre Daten dateibasiert (im \ac{json}-Format) ablegt.

\section{Seiten\"ubersicht}
Die Anwendung besteht aus mehreren Inhaltsseiten mit gemeinsamem Kopf- und Fu\ss{}bereich:
\begin{itemize}
  \item \textbf{Startseite} (\texttt{index.php}) -- Einstieg mit Hero-Bereich und Kennzahlen.
  \item \textbf{\"Uber uns} (\texttt{about.php}) -- Vorstellung des (fiktiven) Teams und der
        Unternehmenswerte.
  \item \textbf{Touren / Kundenstimmen} (\texttt{tours.php}) -- Tour\-angebote und mehrere
        Kundenbewertungen.
  \item \textbf{Buchung} (\texttt{booking.php}) -- Buchungsformular mit Echtzeit-Kostenanzeige.
  \item \textbf{Best\"atigung} (\texttt{confirmation.php}) -- Zusammenfassung nach erfolgreicher
        Buchung.
  \item \textbf{Rechtliches} (\texttt{legal.php}) -- Impressum, Datenschutz und AGB.
  \item \textbf{Formular-Test} (\texttt{formcheck.php}) -- kleine Seite zur Demonstration der
        GET-/POST-Verarbeitung.
\end{itemize}

\section{Gestaltung und technischer Umfang}
Das Layout folgt der Struktur Kopfbereich--Inhalt--Fu\ss{}bereich und ist responsiv umgesetzt.
Die Farbgebung (Gelb-, Orange- und Gr\"unt\"one) ist \"uber \ac{css}-Variablen zentral definiert,
wodurch sich auch ein Hell-/Dunkel-Modus realisieren l\"asst. Das Logo ist vollst\"andig in
\ac{css} gestaltet (ohne Bilddatei). Die auf den Seiten verwendeten St\"adtebilder wurden mit
einem KI-Bildgenerator erzeugt (vgl.\ Kapitel~\ref{ch:b-kitools}).

Der technische Umfang umfasst sieben PHP-Seiten, zwei gemeinsame Includes (Kopf- und
Fu\ss{}bereich), zwei JavaScript-Dateien (Buchungslogik und allgemeine Oberfl\"achenfunktionen)
sowie ein zentrales Stylesheet. Eine Datenbank wird nicht ben\"otigt; Buchungen werden als
\ac{json}-Dateien gespeichert.
